% !TeX root=../main.tex
\chapter{مرور ادبیات}

\section{مفاهیم و تعاریف}
در این بخش، به تبیین مفاهیم پایه، تعاریف فنی و مبانی نظری پرداخته می‌شود که برای درک دقیق مسئله تخمین عمق و نوآوری‌های مطرح‌شده در این پژوهش ضروری هستند.

\subsection{تخمین عمق تک‌دوربینه}
تخمین عمق تک‌دوربینه\enfootnote{Monocular Depth Estimation} به فرآیند استخراج اطلاعات سه‌بعدی و فاصله نقاط صحنه از مرکز نوری دوربین، تنها با استفاده از یک تصویر دوبعدی منفرد اطلاق می‌شود~\cite{Eigen2014Depth}. برخلاف سیستم‌های استریو که از اختلاف منظر میان دو تصویر برای محاسبه عمق استفاده می‌کنند، این رویکرد به دلیل اتکا به یک نمای واحد، ذاتاً یک مسئله «بد‌حال»\enfootnote{Ill-posed Problem} محسوب می‌شود؛ زیرا بی‌نهایت صحنه سه‌بعدی متفاوت می‌توانند منجر به ایجاد یک تصویر دوبعدی یکسان شوند~\cite{Koch1998IllPosed}. مدل‌های مدرن در این حوزه، از ویژگی‌های سطح بالا نظیر پرسپکتیو، بافت و اندازه نسبی اشیا برای غلبه بر این ابهام هندسی بهره می‌برند.

\subsection{عمق متریک در مقابل عمق غیرمتریک}
یکی از چالش‌های بنیادین در بینایی ماشین، تمایز میان عمق نسبی (غیرمتریک\enfootnote{Non-metric Depth}) و عمق مطلق (متریک\enfootnote{Metric Depth}) است. در تخمین عمق غیرمتریک، هدف بازیابی ترتیب قرارگیری اشیا و ساختار کلی صحنه است، بدون آنکه فواصل با واحدهای فیزیکی (مانند متر) مشخص باشند. در مقابل، عمق متریک به معنای بازیابی فواصل واقعی فیزیکی است که برای کاربردهای عملیاتی نظیر ناوبری رباتیک و سیستم‌های کمک‌راننده حیاتی است~\cite{Ranftl2021DPT}. بسیاری از مدل‌های بزرگ مقیاس فعلی به دلیل آموزش بر روی داده‌های متنوع و ناهمگن، خروجی‌های غیرمتریک تولید می‌کنند که فاقد مقیاس فیزیکی پایدار هستند.

\subsection{ابهام مقیاس}
ابهام مقیاس\enfootnote{Scale Ambiguity} به ناتوانی مدل در تشخیص اندازه واقعی اشیا و فاصله آن‌ها از دوربین بدون داشتن اطلاعات اضافی اطلاق می‌شود. در تصویربرداری تک‌دوربینه، یک جسم کوچک در فاصله نزدیک و یک جسم بزرگ در فاصله دور می‌توانند تصویری مشابه روی حسگر ایجاد کنند. این پدیده باعث می‌شود که مدل‌های یادگیری عمیق در بازیابی مقیاس فیزیکی دچار خطا شوند، مگر آنکه اطلاعاتی درباره هندسه دوربین یا ابعاد شناخته‌شده در صحنه در اختیار داشته باشند~\cite{Zhou2017Unsupervised}.

\subsection{پارامترهای درونی دوربین}
پارامترهای درونی دوربین\enfootnote{Camera Intrinsics} مجموعه‌ای از ویژگی‌های هندسی و اپتیکی هستند که نحوه نگاشت نقاط از فضای سه‌بعدی به صفحه تصویر را تعیین می‌کنند. این پارامترها معمولاً در قالب یک ماتریس مشخصه\enfootnote{Intrinsic Matrix} (ماتریس $K$) تعریف می‌شوند که شامل فاصله کانونی\enfootnote{Focal Length}، مرکز تصویر\enfootnote{Principal Point} و ضریب پیچش است~\cite{HartleyZisserman2004}. ادغام صریح این پارامترها در مدل‌های یادگیری عمیق، امکان تبدیل اطلاعات تصویری به مختصات هندسی واقعی را فراهم کرده و نقشی کلیدی در گذار از تخمین عمق نسبی به متریک ایفا می‌کند.

\subsection{ترنسفورمر بینایی}
ترنسفورمر بینایی\enfootnote{Vision Transformer (ViT)} معماری پیشرفته‌ای در یادگیری عمیق است که از مکانیسم‌های خود‌توجهی\enfootnote{Self-attention} برای پردازش تصاویر استفاده می‌کند~\cite{Dosovitskiy2020ViT}. برخلاف شبکه‌های عصبی کانولوشنی که بر ویژگی‌های محلی تمرکز دارند، ترنسفورمرها قادر به مدل‌سازی وابستگی‌های طولانی‌دامنه در کل تصویر هستند~\cite{Ranftl2021DPT}. در مدل پایه این پژوهش، یعنی \lr{Depth Anything V2}، از این معماری برای استخراج ویژگی‌های غنی و مقاوم در برابر تغییرات محیطی استفاده شده است که منجر به تعمیم‌پذیری بالای مدل در سناریوهای صفر‌نمونه\enfootnote{Zero-shot Generalization} می‌شود~\cite{DepthAnythingV2}.

\subsection{تقطیر دانش}
تقطیر دانش\enfootnote{Knowledge Distillation} یک راهبرد آموزشی است که در آن یک مدل بزرگ و پیچیده (مدل معلم\enfootnote{Teacher Model})، دانش خود را به یک مدل دیگر (مدل شاگرد\enfootnote{Student Model}) منتقل می‌کند~\cite{Hinton2015Distillation}. در این پژوهش، از این رویکرد برای انتقال دانش ساختاری و ادراکی از یک مدل تخمین عمق غیرمتریک قدرتمند به مدلی استفاده شده است که با هدف یادگیری مقیاس فیزیکی و ادغام اطلاعات هندسی طراحی شده است. این فرآیند باعث می‌شود مدل شاگرد علاوه بر حفظ دقت در جزئیات ساختاری، پایداری بیشتری در تخمین مقیاس متریک داشته باشد.

\begin{figure}[htbp]
    \centering
    \begin{tikzpicture}[
        block/.style={draw, thick, rectangle, minimum width=2.5cm, minimum height=1cm, fill=orange!10, rounded corners},
        arrow/.style={thick, -Latex}
    ]
    
    \node[block, fill=blue!10] (teacher) at (0, 2) {مدل معلم (غیرمتریک)};
    \node[block, fill=green!10] (student) at (0, -1.5) {مدل شاگرد (متریک)};
    
    \node[draw, dashed, circle] (image) at (-4, 0.25) {تصویر ورودی};
    
    \draw[arrow] (image) |- (teacher);
    \draw[arrow] (image) |- (student);
    
    \node[draw, rectangle, rounded corners] (t_out) at (4, 2) {نقشه عمق معلم};
    \node[draw, rectangle, rounded corners] (s_out) at (4, -1.5) {نقشه عمق شاگرد};
    
    \draw[arrow] (teacher) -- (t_out);
    \draw[arrow] (student) -- (s_out);
    
    \draw[arrow, dashed, red, thick] (t_out) -- (s_out) node[midway, right] {انتقال دانش (تطبیق)};
    
    \end{tikzpicture}
    \caption{مفهوم کلی تقطیر دانش در تخمین عمق: مدل معلم دانش ساختاری و ویژگی‌های بصری خود را به مدل شاگرد منتقل می‌کند تا مدل شاگرد ضمن یادگیری مقیاس فیزیکی، دقت ساختاری خود را حفظ کند.}
    \label{fig:knowledge_distillation}
\end{figure}

\subsection{یادگیری خود‌نظارتی}
یادگیری خود‌نظارتی\enfootnote{Self-supervised Learning} روشی از یادگیری است که در آن مدل بدون نیاز به برچسب‌های دستی (نقشه‌های عمق مرجع)، و تنها با استفاده از ساختار درونی داده‌ها آموزش می‌بیند. در تخمین عمق، این امر معمولاً از طریق بازسازی دیدگاه\enfootnote{View Synthesis} و استفاده از هندسه اپی‌پولار در ویدئوها یا تصاویر استریو محقق می‌شود~\cite{Godard2017Unsupervised}. این رویکرد امکان استفاده از حجم عظیمی از داده‌های بدون برچسب را فراهم کرده و وابستگی به سنسورهای گران‌قیمت نظیر لیدار را کاهش می‌دهد.

\section{مرور ادبیات و پیشینه پژوهش}
پژوهش‌های انجام شده در حوزه تخمین عمق تک‌دوربینه در طول دو دهه اخیر، گذار از مدل‌های آماری ساده به شبکه‌های عصبی بسیار پیچیده را تجربه کرده‌اند. در این بخش، این مطالعات بر اساس ماهیت روش‌شناختی و نحوه مواجهه با چالش‌های هندسی دسته‌بندی و تحلیل می‌شوند.

\begin{figure}[htbp]
    \centering
    \begin{tikzpicture}[
        every node/.style={draw, thick, align=center, fill=blue!10, rounded corners},
        level 1/.style={sibling distance=5cm, level distance=1.5cm},
        level 2/.style={sibling distance=2.5cm, level distance=1.5cm},
        edge from parent/.style={draw, thick, -Latex}
    ]
    \node {روش‌های تخمین عمق تک‌دوربینه}
        child {node {کلاسیک و مبتنی\\بر ویژگی دستی}
            child {node {میدان‌های تصادفی\\مارکوف (MRF)}}
        }
        child {node {یادگیری عمیق}
            child {node {نظارت‌شده}}
            child {node {خودنظارتی}
                child {node {مبتنی بر استریو}}
                child {node {مبتنی بر ویدئو}}
            }
        };
    \end{tikzpicture}
    \caption{دسته‌بندی کلی روش‌های تخمین عمق تک‌دوربینه در ادبیات موضوع.}
    \label{fig:mde_taxonomy}
\end{figure}

\subsection{رویکردهای کلاسیک و یادگیری مبتنی بر ویژگی‌های دستی}
تلاش‌های اولیه برای استخراج عمق از تصاویر منفرد، عمدتاً بر نشانه‌های پرسپکتیو\enfootnote{Perspective Cues} و ویژگی‌های هندسی صریح متکی بودند. یکی از کارهای پیشگام در این حوزه، سیستم \lr{Make3D} توسط \cite{Saxena2009Make3D} ارائه شد که از یک مدل میدان تصادفی مارکوف\enfootnote{Markov Random Field (MRF)} برای مدل‌سازی روابط مکانی میان قطعات تصویر استفاده می‌کرد~\cite{Saxena2009Make3D}. اگرچه این روش‌ها در بازسازی ساختارهای ساده موفق بودند، اما به دلیل وابستگی شدید به ویژگی‌های دستی\enfootnote{Hand-crafted Features}، در محیط‌های پیچیده و متنوع با شکست مواجه می‌شدند. در واقع، ناتوانی این مدل‌ها در استخراج مفاهیم معنایی سطح بالا از تصویر، منجر به محدودیت شدید در تعمیم‌پذیری آن‌ها می‌گردید.

\subsection{تحول در تخمین عمق با یادگیری عمیق نظارت‌شده}
با ظهور شبکه‌های عصبی عمیق، رویکرد یادگیری نظارت‌شده به استاندارد طلایی این حوزه تبدیل شد. \cite{Eigen2014Depth} برای نخستین بار نشان دادند که استفاده از یک معماری دو مقیاسه شامل شبکه‌های کانولوشنی می‌تواند ساختار کلی و جزئیات محلی عمق را به‌طور هم‌زمان بیاموزد~\cite{Eigen2014Depth}. در ادامه، این تیم تحقیقاتی با ترکیب ویژگی‌های معنایی و عمق، دقت مدل را بهبود بخشیدند~\cite{Eigen2015Predicting}. با این حال، همان‌طور که در پژوهش‌های بعدی اشاره شد، این مدل‌ها به دلیل آموزش بر روی دیتاست‌های محدود نظیر \lr{NYU Depth V2}~\cite{Silberman2012NYUv2}، در مواجهه با داده‌های خارج از توزیع\enfootnote{Out-of-distribution Data} دچار افت عملکرد جدی می‌شدند و خروجی آن‌ها به‌شدت به بافت تصویر وابسته بود تا هندسه واقعی.

\subsection{یادگیری خودنظارتی و بازسازی هندسی دیدگاه}
به منظور کاهش وابستگی به داده‌های برچسب‌دار لیدار، محققان به سمت یادگیری خودنظارتی متمایل شدند. \cite{Godard2017Unsupervised} با استفاده از انسجام زمانی در ویدئوها و روابط استریو، مسئله تخمین عمق را به یک مسئله بازسازی تصویر\enfootnote{Image Reconstruction} تبدیل کردند~\cite{Godard2017Unsupervised}. کار آن‌ها در مدل \lr{Monodepth2} با معرفی توابع هزینه مقاوم در برابر انسداد و استفاده از شبکه‌های تخمین ژست دوربین\enfootnote{Pose Estimation Network} به تکامل رسید~\cite{Godard2019Digging}. با وجود این پیشرفت‌ها، یک چالش اساسی همچنان باقی ماند: ابهام مقیاس. در این روش‌ها، مدل تنها قادر به بازیابی عمق تا یک ضریب مقیاس مجهول بود که مانع از کاربرد مستقیم آن‌ها در سیستم‌های ناوبری متریک می‌شد.

\subsection{مدل‌های مبتنی بر ترنسفورمر و یادگیری در مقیاس بزرگ}
در سال‌های اخیر، استفاده از معماری ترنسفورمر بینایی فصلی نو در تخمین عمق گشوده است. مدل‌هایی مانند \lr{DPT} نشان دادند که جایگزینی لایه‌های کانولوشنی با مکانیسم‌های توجه جهانی می‌تواند به بازسازی مرزهای دقیق‌تر و درک بهتر از ساختار جهانی صحنه منجر شود~\cite{Ranftl2021DPT}. اخیراً مدل \lr{Depth Anything V2} با بهره‌گیری از داده‌های عظیم و تکنیک‌های پیشرفته تقطیر دانش، به پایداری بی‌سابقه‌ای در تخمین عمق غیرمتریک دست یافته است~\cite{DepthAnythingV2}. با این حال، این مدل‌های قدرتمند نیز همچنان اطلاعات درونی دوربین را نادیده می‌گیرند. در حالی که \cite{Ranftl2021DPT} بر تعمیم‌پذیری در محیط‌های مختلف تأید داشتند، نادیده گرفتن پارامترهای هندسی باعث می‌شود که خروجی این مدل‌ها با تغییر فاصله کانونی دوربین دچار نوسانات مقیاسی شود.

\subsection{ادغام دانش هندسی و چالش پایداری متریک}
تلاش‌های محدودی برای پیوند میان هندسه کلاسیک و یادگیری عمیق صورت گرفته است. \cite{Zhou2017Unsupervised} نشان دادند که بدون داشتن مقیاس پایه، تخمین عمق مطلق غیرممکن است~\cite{Zhou2017Unsupervised}. برخی پژوهش‌ها نظیر کار \cite{Guizilini2020Semantically} تلاش کردند با استفاده از ویژگی‌های معنایی هدایت‌شده توسط هندسه، مقیاس را بازیابی کنند~\cite{Guizilini2020Semantically}. با این حال، خلأ موجود در ادبیات فعلی، نبود یک راهکار نظام‌مند برای تزریق مستقیم ماتریس مشخصه دوربین به دیکودرهای مدل‌های بزرگ‌مقیاس است. اکثر روش‌های فعلی یا بر کالیبراسیون پس‌پردازشی تکیه دارند و یا فرض می‌کنند شبکه به‌صورت ضمنی هندسه را می‌آموزد؛ فرضی که در شرایط تغییر دوربین به چالش کشیده می‌شود. پژوهش حاضر بر ادغام صریح پارامترهای درونی در معماری \lr{Depth Anything V2} تمرکز دارد.

\section{شکاف تحقیقاتی}
با وجود پیشرفت‌های چشمگیر در حوزه تخمین عمق تک‌دوربینه، تحلیل انتقادی ادبیات موضوع نشان‌دهنده وجود خلأهای تحقیقاتی جدی است که مانع از کاربرد گسترده این فناوری در محیط‌های عملیاتی می‌شود. 

\subsection{فقدان پایداری در تخمین عمق متریک}
بخش عمده‌ای از پژوهش‌های اخیر بر بهبود دقت نسبی معطوف شده‌اند~\cite{DepthAnythingV2, Ranftl2021DPT}. با این حال، این مدل‌ها در تولید خروجی با مقیاس فیزیکی واقعی دچار ضعف هستند. شکاف موجود، ناتوانی مدل‌ها در حفظ پایداری مقیاس\enfootnote{Scale Stability} در مواجهه با صحنه‌هایی است که از نظر بصری مشابه اما از نظر هندسی متفاوت هستند. فقدان یک مکانیسم درونی برای تبدیل ویژگی‌های بصری به مقادیر فیزیکی، باعث می‌شود که این مدل‌ها بدون کالیبراسیون خارجی در کاربردهای متریک غیرقابل استفاده باقی بمانند.

\subsection{نادیده گرفتن صریح هندسه و پارامترهای درونی دوربین}
اکثر مدل‌های فعلی فرض می‌کنند که شبکه عصبی می‌تواند پارامترهای درونی دوربین را به‌صورت ضمنی بیاموزد. با این حال، تغییر در مشخصات سنسور یا لنز منجر به بروز خطا در تخمین عمق می‌شود~\cite{Zhou2017Unsupervised}. خلأ اصلی تحقیقاتی، نبودِ یک معماری یکپارچه است که بتواند ماتریس مشخصه دوربین را به‌عنوان یک ورودی مرتبه اول پردازش کند. اکثر روش‌ها یا هندسه را نادیده می‌گیرند و یا از آن تنها در مراحل پس‌پردازش استفاده می‌کنند که منجر به عدم انطباق ویژگی‌ها با قیود هندسی در لایه‌های میانی شبکه می‌شود.

\subsection{حساسیت بالا نسبت به تغییرات دامنه و داده‌های خارج از توزیع}
مدل‌هایی که بر روی دیتاست‌های خاص آموزش دیده‌اند، در مواجهه با داده‌های خارج از توزیع دچار پدیده لغزش مقیاس\enfootnote{Scale Drift} می‌شوند~\cite{Ranftl2019Robust}. اگرچه روش‌های تقطیر دانش تلاش کرده‌اند این مشکل را حل کنند~\cite{DepthAnythingV2}، اما وابستگی خروجی متریک به توزیع داده‌های آموزشی همچنان یک چالش است. شکاف تحقیقاتی در این است که پایداری متریک باید بر اساس درک صریح از رابطه فیزیکی پیکسل‌ها و دوربین باشد، نه صرفاً حافظه آماری از داده‌ها.

\subsection{محدودیت معیارهای ارزیابی در تفکیک خطاهای هندسی}
در بسیاری از مطالعات، از معیارهای نرمال‌شده‌ای نظیر \lr{AbsRel} استفاده می‌شود که قادر به آشکارسازی خطاهای ناشی از عدم پایداری مقیاس فیزیکی نیستند~\cite{Ranftl2021DPT}. خلأ موجود، کمبود تحلیل‌های عمیق بر روی معیارهای حساس به واحد فیزیکی (مانند \lr{RMSE}) در سناریوهایی است که پارامترهای درونی دوربین تغییر می‌کنند.

\section{نوآوری‌های تحقیق}
پژوهش حاضر با هدف گذار به تخمین عمق هندسه‌آگاه\enfootnote{Geometry-aware Depth Estimation}، چندین نوآوری کلیدی را ارائه می‌دهد:

\subsection{ادغام صریح پارامترهای درونی دوربین در معماری ترنسفورمر}
نوآوری اصلی این پژوهش، تزریق مستقیم و صریح ماتریس مشخصه دوربین به درون ساختار دیکودر مدل \lr{Depth Anything V2} است. در این رویکرد، پارامترهای هندسی پس از نرمال‌سازی و نگاشت به یک فضای نهان\enfootnote{Latent Space}، به‌عنوان اطلاعات مکمل در کنار ویژگی‌های بصری قرار می‌گیرند. این نوآوری مستقیماً شکاف مربوط به ناپایداری مقیاس را هدف قرار می‌دهد.

\subsection{طراحی دیکودر چند‌مؤلفه‌ای با قابلیت تفکیک ویژگی‌های بصری و هندسی}
در این پژوهش، معماری دیکودر به‌گونه‌ای اصلاح شده است که با استفاده از مکانیزم‌های ادغام ویژگی\enfootnote{Feature Fusion Mechanisms}، وزن دهی به ویژگی‌های بصری را بر اساس پارامترهای درونی دوربین تنظیم کند. این ساختار اجازه می‌دهد که مدل در سناریوهای صفرنمونه، پایداری متریک خود را حتی در صورت تغییرات شدید در فاصله کانونی حفظ نماید.

\subsection{راهبرد آموزش چند‌مرحله‌ای مبتنی بر تقطیر دانش متریک}
این تحقیق یک چارچوب آموزشی جدید پیشنهاد می‌دهد که در آن از تقطیر دانش~\cite{Hinton2015Distillation} برای انتقال پایداریِ ادراکی از یک مدل معلم غیرمتریک به یک مدل شاگرد متریک استفاده می‌شود. نوآوری در طراحی تابع هزینه ترکیبی است که هم‌زمان بر حفظ ساختار لبه‌ها و انطباق مقیاس فیزیکی با داده‌های واقعیت‌میدانی\enfootnote{Ground Truth} تمرکز دارد.

\subsection{تحلیل حساسیت سیستماتیک نسبت به تغییرات هندسه تصویربرداری}
این پژوهش با ارائه یک چارچوب ارزیابی دقیق، تأثیر تغییر پارامترهای دوربین را بر معیارهای مختلف بررسی می‌کند. این نوآوری تحلیلی نشان می‌دهد که چگونه معیارهای سنتی ممکن است در نمایش خطاهای مقیاس ناتوان باشند و در مقابل، معیارهایی نظیر \lr{RMSE} و \lr{SiLog} رفتار هندسی واقعی مدل را منعکس می‌کنند~\cite{Godard2019Digging}.

\section{جمع‌بندی فصل}
در این فصل، با بررسی پیشینه پژوهش مشخص گردید که مدل‌های فعلی علی‌رغم دقت بصری بالا، در حفظ «پایداری مقیاس متریک» دچار ضعف هستند. این چالش ریشه در نادیده گرفتن صریح پارامترهای هندسی دوربین دارد که منجر به ابهام مقیاس در محیط‌های متنوع می‌شود. پژوهش حاضر با هدف رفع این شکاف، بر ادغام مستقیم ماتریس مشخصه دوربین در معماری مدل \lr{Depth Anything V2} تمرکز دارد. در فصل آینده، جزئیات متدولوژی، فرمول‌بندی ریاضی این ادغام و راهبرد آموزش مدل پیشنهادی به‌صورت تفصیلی تشریح خواهد شد.